%!TeX root=main-DS1.tex
\chapter{Binarne relacije}

\v Zivot u savremenom svetu je postao nezamisliv bez procesa koji podrazumevaju obradu \emph{ogromnih} koli\v cina podataka.
Zato i ne \v cudi \v sto se pojavljuju nove nau\v cne discipline (kao \v sto je nauka o podacima -- \emph{data science})
koje su posve\'cene upravo tome.
Mi, ljudi, kao vrsta imamo ozbiljna mentalna ograni\v cenja kada se radi o mogu\'cnosti
sagledavanja i analize velikih koli\v cina podataka. Da bismo mogli da shvatimo savremene fenomene u nauci ili dru\v stvu
moramo da se dovijamo na razne na\v cine kako bismo se izborili sa sve ve\'cim nivoom slo\v zenosti.

Dve osnove strategije koje koristimo da se emotivno izborimo sa velikim koli\v cinama podataka su
klasifikacija (grupisanje) podataka, i sortiranje (ure\dj ivanje) podataka.
Klasifikacija podataka se svodi na to da se podaci organizuju u grupe (klase) podataka
koji su \enquote{sli\v cni} u odnosu na neki kriterijum. S druge strane, sortiranje podataka se svodi na to da se
podaci pore\dj aju u odnosu na neki kriterijum od \enquote{manjih} vrednosti ka \enquote{ve\'cim} ili obrnuto.

U oba slu\v caja imamo potrebu da iz nekog skupa podataka $A$ uzmemo dva podatka
$a, b \in A$ i da utvrdimo da li su oni u vezi: da li je $a$ \enquote{sli\v can} sa $b$ u prvom slu\v caju, odnosno
da li je $a$ \enquote{manji} od $b$ u drugom slu\v caju.
U apstraktnom smislu, to zna\v ci da je potrebno da uspostavimo binarnu relaciju
$\rho \subseteq A \times A$ na skupu $A$ takvu da
$(a, b) \in \rho$ modeluje situaciju \enquote{$a$ je sli\v can sa $b$} u prvom slu\v caju,
odnosno, \enquote{$a$ je manji od $b$} u drugom slu\v caju.

\section{Klasifikacija}

U ovom odeljku \'cemo odgovoriti na pitanje
kada neka apstraktna binarna relacija $\rho \subseteq A \times A$ predstavlja
\enquote{razuman} kriterijum za klasifikaciju.

\begin{figure}
    \centering
    \input ft07-slicnost.pgf
    \caption{Sli\v cnost trouglova}    
    \label{ft07.fig.slicnost}
\end{figure}

\begin{EX}
Za dva trougla ka\v zemo da su \emph{sli\v cni} ako imaju iste
uglove. Na primer, trouglovi $T$ i $P$ na Sl.~\ref{ft07.fig.slicnost} nisu sli\v cni, \v sto
pi\v semo $T \not\sim P$, dok su trouglovi $P$, $Q$ i $R$ me\dj usobno sli\v cni,
$P \sim Q$, $P \sim R$, $Q \sim R$.
Lako je videti da relacija sli\v cnosti $\sim$ ima slede\'ce osobine:
\begin{itemize}\itemsep -1.5pt
\item svaki trougao je sli\v can sebi: $X \sim X$,
\item ako je trougao $X$ sli\v can trouglu $Y$ onda je i trougao $Y$ sli\v can trouglu $X$:\newline
      $X \sim Y \Rightarrow Y \sim X$, i
\item ako su trouglovi $X$ i $Y$ sli\v cni, i ako su touglovi $Y$ i $Z$ sli\v cni, onda su i trouglovi $X$ i $Z$ sli\v cni:
      $X \sim Y \land Y \sim Z \Rightarrow X \sim Z$.
\end{itemize}
\end{EX}

\begin{EX}
    Relacija kongruentnosti $\equiv_m$ ima osobine navedene u prethodnom primeru:
    \begin{itemize}
    \item za svako $a \in \NN$ va\v zi: $a \equiv_m a$,
    \item za sve $a, b \in \NN$ va\v zi: $a \equiv_m b \Rightarrow b \equiv_m a$, i
    \item za sve $a, b, c \in \NN$: $a \equiv_m b \land b \equiv_m c \Rightarrow a \equiv_m c$.
    \end{itemize}
\end{EX}

Ove tri osobine karakteri\v su pojam \enquote{sli\v cnosti} objekata s obzirom na neko svojstvo, a uop\v stavanjem
dobijamo pojam \emph{relacije ekvivalencije} koji predstavlja apstrakciju \enquote{razumnog} kriterijuma za klasifikaciju podataka.

\begin{DEF}
    \emph{Relacija ekvivalencije (\enquote{razuman} kriterijum za klasifikaciju) na skupu $A$} je svaka binarna relacija $\rho \subseteq A \times A$
    koja ima slede\'ce tri osobine:
    \begin{itemize}\itemsep -1pt
    \item za sve $a \in A$: $a \mathrel{\rho} a$, \hfill (\emph{refleksivnost})
    \item za sve $a, b \in A$: $a \mathrel{\rho} b \Rightarrow b \mathrel{\rho} a$, \hfill (\emph{simetri\v cnost})
    \item za sve $a, b, c \in A$: $a \mathrel{\rho} b \land b \mathrel{\rho} c \Rightarrow a \mathrel{\rho} c$. \hfill (\emph{tranzitivnost})
    \end{itemize}
\end{DEF}

\begin{EX}\label{ft07.ex.01-weight}
  Neka je $n \ge 1$ prirodan broj i neka je $A = \{0,1\}^n$ skup svih 01-nizova du\v zine $n$. Na skupu $A$ defini\v semo
  relaciju $\sim$ ovako: za 01-nizove $u$ i $v$ stavimo $u \sim v$ ako i samo ako nizovi $u$ i $v$ imaju isti broj jedinica
  (broj jedinica u 01-nizu se zove \emph{te\v zina} tog niza; dakle dva niza su ekvivalentni ako imaju istu te\v zinu).
  Na primer, za $n = 8$ je $00001111 \sim 01010110$ i $00110011 \not\sim 11001000$.
  Lako se vidi da je ovo relacija ekvivalencije na $A$.
\end{EX}

\begin{EX}\label{ft07.ex.iskazne-fle}
  Neka je $\Phi$ skup svih iskaznih formula. Na skupu $\Phi$ smo definisali relaciju $\equiv$ ovako:
  $\phi \equiv \psi$ ako i samo ako $\vDash \phi \Leftrightarrow \psi$. Relacija $\equiv$ je relacija
  ekvivalencije na $\Phi$ zbog slede\'ce tri tautologije:
  \begin{itemize}\itemsep -1.5pt
  \item
    Refleksivnost: $\vDash \phi \Leftrightarrow \phi$ za svaku iskaznu formulu $\phi$;
  \item
    Simetri\v cnost: $\vDash (\phi \Leftrightarrow \psi) \Leftrightarrow (\psi \Leftrightarrow \phi)$;
  \item
    Tranzitivnost: $\vDash (\phi \Leftrightarrow \psi) \land (\psi \Leftrightarrow \theta) \Rightarrow (\phi \Leftrightarrow \theta)$.
  \end{itemize}
\end{EX}

Sada \'cemo pokazati da se relacija ekvivalencije pona\v sa kao kriterijum za klasifikaciju
podataka tako \v sto \'cemo pokazati da svaka takva relacija razbija skup na kome je definisana
na klastere \enquote{me\dj usobno sli\v cnih} objekata, \v sto i jeste cilj svake klasifikacije.

\begin{figure}
  \centering
  \input ft07-01-nizovi.pgf
  \caption{Relacija $\sim$ iz Primera~\ref{ft07.ex.01-weight-2}}
  \label{ds1.sl.01-4}
\end{figure}

\begin{EX}\label{ft07.ex.01-weight-2}
  Posmatrajmo ponovo situaciju opisanu u Primeru~\ref{ft07.ex.01-weight} s tim da \'cemo radi odre\dj enosti
  uzeti da je $A = \{0,1\}^4$. Podsetimo se da za 01-nizove $u$ i $v$ sa $u \sim v$ ozna\v cavamo
  da $u$ i $v$ imaju istu te\v zinu (= isti broj jedinica). Relacija $\sim$ je ilustrovana na Sl.~\ref{ds1.sl.01-4}.
  Vidimo da su elementi skupa $A$ podeljeni u grupe (klastere), i da se u istoj grupi nalaze
  svi elementi koji su me\dj usobno \enquote{sli\v cni}, \v sto u ovom slu\v caju zna\v ci da imaju istu
  te\v zinu. Niz 0000 je sam u klasi jer je to jedini niz sa te\v zinom~0. Sli\v cno je i niz 1111
  sam u klasi. Relacija $\sim$ je tako klasifikovala elemente skupa $A$ po te\v zini.
\end{EX}

\begin{DEF}
  Neka je $\Boxed\rho \subseteq A \times A$ relacija ekvivalencije na $A$ (\emph{kriterijum za klasifikaciju}).
  \emph{Klaster} elementa $a \in A$ je skup svih elemenata skupa $A$ koji su \enquote{sli\v cni} sa $a$
  u odnosu na svojstvo modelovano relacijom $\rho$:
  $$
    [a]_\rho = \{ b \in A : a \mathrel{\rho} b \} = \{ b \in A : b \mathrel{\rho} a \}.
  $$
  (Poslednja jednakost je posledica simetri\v cnosti relacije $\rho$.)
\end{DEF}

\begin{EX}
  Evo klastera nekih od elemenata skupa $A$ iz Primera~\ref{ft07.ex.01-weight-2}:
  \begin{align*}
    [0000]_\sim &= \{0000\}\\
    [0001]_\sim &= \{0001, 0010, 0100, 1000\}\\
    [0010]_\sim &= \{0001, 0010, 0100, 1000\}\\
    [0100]_\sim &= \{0001, 0010, 0100, 1000\}\\
    [0011]_\sim &= \{0011, 0101, 1001, 0110, 1010, 1100\}\\
    [1100]_\sim &= \{0011, 0101, 1001, 0110, 1010, 1100\}.
  \end{align*}
  Prime\'cujemo: klaster elementa $a$ sadr\v zi element $a$ (naravno!), i mo\v zda jo\v s neke elemente;
  pored toga, ako je $u \sim v$, onda $u$ i $v$ pripadaju istom klasteru.
\end{EX}

Sada \'cemo sve \v sto smo uo\v cili na ovom malom primeru pokazati i formalno.

\begin{THM}\label{ft07.thm.ekviv}
  Neka je $\Boxed\rho \subseteq A \times A$ relacija ekvivalencije na $A$ (kriterijum za klasifikaciju).
  \begin{enumerate}\itemsep -1pt
  \item
    Svaki element skupa $A$ pripada svom klasteru: $a \in [a]_\rho$.
  \item
    Elementi koji su u relaciji odre\dj uju isti klaster:
    $a \mathrel{\rho} b \Rightarrow [a]_\rho = [b]_\rho$.
  \item
    Razli\v citi klasteri su disjunktni:
    $[a]_\rho \ne [b]_\rho \Rightarrow [a]_\rho \cap [b]_\rho = \0$.
  \item
    Neka je $C_1$, \ldots, $C_k$ spisak svih klastera relacije $\rho$. Tada je
	$A = C_1 \cup \ldots \cup C_k$.
  \end{enumerate}
\end{THM}
\begin{proof}[Dokaz]
  (1)
  Neka je $a \in A$ proizvoljno. Zbog refleksivnosti je $a \mathrel{\rho} a$, pa je $a \in [a]_\rho$.
  
  \medskip
  
  (2)
  Neka je $a \mathrel\rho b$. Imaju\'ci u vidu da su klasteri skupovi, da bismo pokazali da je $[a]_\rho = [b]_\rho$,
  potrebno je pokazati da je $[a]_\rho \subseteq [b]_\rho$ i $[b]_\rho \subseteq [a]_\rho$. Pokaza\'cemo samo
  $[a]_\rho \subseteq [b]_\rho$ zato \v sto se $[b]_\rho \subseteq [a]_\rho$ pokazuje na isti na\v cin.
  
  Uzmimo neko $x \in [a]_\rho$. Tada je $x \mathrel\rho a$ prema definiciji klastera. S druge strane, mi smo pretpostavili
  da je $a \mathrel\rho b$, pa nam tranzitivnost relacije $\rho$ daje $x \mathrel\rho b$. Zato je $x \in [b]_\rho$.
  Ovim je pokazano da je $[a]_\rho \subseteq [b]_\rho$.
  
  \medskip
  
  (3)
  Neka su $a, b \in A$. Dokazujemo kontrapoziciju tvr\dj enja:
  $$
    [a]_\rho \cap [b]_\rho \ne \0 \Rightarrow  [a]_\rho = [b]_\rho.
  $$
  Neka $[a]_\rho \cap [b]_\rho \ne \0$ i neka je $x \in [a]_\rho \cap [b]_\rho$ proizvoljno.
  Prema definiciji klastera, $x \in [a]_\rho$ zna\v ci da je $x \mathrel\rho a$, pa na osnovu (2) zaklju\v cujemo da je $[x]_\rho = [a]_\rho$.
  Sli\v cno, $x \in [b]_\rho$ zna\v ci da je $x \mathrel\rho b$, pa na osnovu (2) zaklju\v cujemo da je $[x]_\rho = [b]_\rho$.
  Dakle, $[a]_\rho = [x]_\rho = [b]_\rho$.

  \medskip
  
  (4)
  Da bismo dokazali da je $A = C_1 \cup \ldots \cup C_k$, dokaza\'cemo da je
  $A \subseteq C_1 \cup \ldots \cup C_k$ i $C_1 \cup \ldots \cup C_k \subseteq A$.

  Uzmimo proizvoljno $a \in A$. Kako smo videli u (1), $a \in [a]_\rho$. No, $[a]_\rho$ je neki od klastera sa spiska, recimo $C_i$.
  Zato je $a \in [a]_\rho = C_i \subseteq C_1 \cup \ldots \cup C_k$.

  S druge strane, svaki klaster je podskup skupa $A$. Prema tome, $C_1 \subseteq A$, \ldots, $C_k \subseteq A$,
  pa onda mora biti i unija svih klastera sadr\v zana u $A$: $C_1 \cup \ldots \cup C_k \subseteq A$.
\end{proof}

\begin{DEF}
  Neka je $\rho \subseteq A \times A$ relacija ekvivalencije na skupu $A$ koja razbija skup $A$ na klastere $C_1, \ldots, C_k$.
  To kra\'ce pi\v semo ovako:
  $$
    A / \rho = \{C_1, \ldots, C_k\}.
  $$
  Skup $A / \rho$ se zove jo\v s i \emph{koli\v cni\v cki skup} relacije $\rho$.
\end{DEF}

\begin{EX}
  Posmatrajmo relaciju $\equiv_5$ na skupu $\ZZ$. Podsetimo se da $a \equiv_5 b$ zna\v ci da
  $a$ i $b$ imaju isti ostatak pri deljenju sa~5. (Sli\v cni su kada se podele sa~5.)
  Tada je
  \begin{align*}
    [0]_{\equiv_5} &= \{\ldots, -10, -5, 0, 5, 10, \ldots\},\\
    [1]_{\equiv_5} &= \{\ldots, -9, -4, 1, 6, 11, \ldots\},\\
    [2]_{\equiv_5} &= \{\ldots, -8, -3, 2, 7, 12, \ldots\},\\
    [3]_{\equiv_5} &= \{\ldots, -7, -2, 3, 8, 13, \ldots\},\\
    [4]_{\equiv_5} &= \{\ldots, -6, -1, 4, 9, 14, \ldots\}.
  \end{align*}
  Po\v sto pri deljenju sa~5 mo\v zemo dobiti samo pet razli\v citih ostataka (0, 1, 2, 3, 4), u ovih pet skupova
  su svrstani svi celi brojevi pa zaklju\v cujemo da je:
  $$
    [0]_{\equiv_5}, [1]_{\equiv_5}, [2]_{\equiv_5}, [3]_{\equiv_5}, [4]_{\equiv_5}
  $$
  spisak svih klastera relacije $\equiv_5$. Pi\v semo:
  $$
    \ZZ / \Boxed{\equiv_5} = \{[0]_{\equiv_5}, [1]_{\equiv_5}, [2]_{\equiv_5}, [3]_{\equiv_5}, [4]_{\equiv_5}\}.
  $$
  Uop\v steno, za $m \ge 2$ je
  $$
    \ZZ / \Boxed{\equiv_m} = \{[0]_{\equiv_m}, [1]_{\equiv_m}, [2]_{\equiv_m}, \ldots, [m-1]_{\equiv_m}\}
  $$
  jer pri deljenju sa $m$ mo\v zemo dobiti samo $m$ razli\v citih ostataka: 0, 1, 2, \ldots, $m-1$.
\end{EX}

\begin{EX}
  U Primeru~\ref{ft07.ex.01-weight} smo na skupu $A = \{0,1\}^n$ svih 01-nizova du\v zine $n$, $n \ge 1$,
  definisali relaciju $\sim$ ovako: $u \sim v$ ako i samo ako nizovi $u$ i $v$ imaju isti broj jedinica.
  Neka je $n = 8$. Primetimo, prvo, da je $00000001 \sim 00000010 \sim 00000100 \sim \ldots \sim 10000000$,
  potom $00000011 \sim 00000101 \sim \ldots \sim 11000000$, kao i da je $00000000$ jedini 01-niz $u \in A^8$
  sa osobinom $u \sim 00000000$. Zato je:
  \begin{align*}
    [00000000]_\sim &= \{00000000\}\\
    [00000001]_\sim &= \{ u \in A^8 : u \text{ ima ta\v cno jednu jedinicu} \},\\
    [00000011]_\sim &= \{ u \in A^8 : u \text{ ima ta\v cno dve jedinice} \},\\
                    &\vdots\\
    [01111111]_\sim &= \{ u \in A^8 : u \text{ ima ta\v cno sedam jedinica} \},\\
    [11111111]_\sim &= \{ 11111111 \}.
  \end{align*}
  Ovim su obuhva\'ceni svi 01-nizovi iz $A^8$, pa zaklju\v cujemo da je
  $$
    A^8 / \Boxed\sim = \{[00000000]_\sim, [00000001]_\sim, [00000011]_\sim, \ldots, [11111111]_\sim\}
  $$
  spisak svih klastera relacije~$\sim$.

  Ovaj primer se mo\v ze uop\v stiti na proizvoljno $n \ge 1$. Lako vidimo da je:
  \begin{align*}
    [000\ldots00]_\sim &= \{000\ldots00\}\\
    [000\ldots01]_\sim &= \{ u \in A^n : u \text{ ima ta\v cno jednu jedinicu} \},\\
                    &\vdots\\
    [011\ldots11]_\sim &= \{ u \in A^n : u \text{ ima ta\v cno $n-1$ jedinica} \},\\
    [111\ldots11]_\sim &= \{ 111\ldots11 \}.
  \end{align*}
  Skup $A^n / \Boxed\sim$ ima $n + 1$ elemenata i to su ta\v cno ovih $n + 1$ klastera.
\end{EX}

\begin{EX}
  U Primeru~\ref{ft07.ex.iskazne-fle} smo na skupu $\Phi$ svih iskaznih formula uo\v cili relaciju
  ekvivalencije $\equiv$ ovako: $\phi \equiv \psi$ ako i samo ako $\vDash \phi \Leftrightarrow \psi$.
  Sada se lako vidi da je \hbox{$[p \lor \lnot p]_\equiv$} skup svih tautologija, dok je
  \hbox{$[p \land \lnot p]_\equiv$} skup svih kontradikcija (\emph{kontradikcija} je iskazna formula koja je neta\v cna
  za sve valuacije).
\end{EX}

Videli smo u primerima da relacija ekvivalencije razbija skup na kome je definisana na klastere.
Interesantno je da va\v zi i obrnuto: ukoliko su elementi skupa $A$ unapred podeljeni u klastere,
uvek se mo\v ze na\'ci relacija ekvivalencije koja ba\v s tako klasifikuje elemente skupa~$A$.

\begin{DEF}
  Neka je $A$ skup i neka su $C_1$, $C_2$, \ldots, $C_k$ neki podskupovi skupa $A$.
  Niz $C_1, C_2, \ldots, C_k$ je \emph{klasterizacija skupa $A$} (drugo ime: \emph{particija skupa $A$})
  ako va\v zi slede\'ce:
\end{DEF}

\noindent
\begin{minipage}{3in}
  \begin{itemize}\itemsep -1.5pt
  \item svi $C_1, C_2, \ldots, C_k$ su neprazni;
  \item razli\v citi elementi niza $C_1, C_2, \ldots, C_k$ su disjunktni; i
  \item $C_1 \cup C_2 \cup \ldots \cup C_k = A$.
  \end{itemize}

  \kern -4mm

  Skupovi $C_1$, $C_2$, \ldots, $C_k$ se zovu \emph{klasteri} (drugo ime: \emph{blokovi particije}).

  \medskip

  \hspace*{1.5em}Dakle, klasteri su neprazni podskupovi skupa $A$ sa osobinom da razli\v citi klasteri nemaju zajedni\v ckih
  elemenata, i da svaki element skupa $A$ pripada nekom klasteru. Odatle sledi da svaki element skupa $A$ pripada
  ta\v cno jednom klasteru.
\end{minipage}
\hfill
\begin{minipage}{2.5in}
  \begin{center}
    \input ft07-r5.pgf
  \end{center}
\end{minipage}

\begin{THM}
  Neka je $C_1, C_2, \ldots, C_k$ klasterizacija skupa $A$. Tada postoji relacija ekvivalencije $\rho$ na skupu $A$
  takva da je $A / \Boxed\rho = \{C_1, C_2, \ldots, C_k\}$.
\end{THM}
\begin{proof}[Dokaz]
  Neka je $C_1, C_2, \ldots, C_k$ klasterizacija skupa $A$ i neka je $\rho$ binarna relacija na $A$ definisana ovako:
  $$
    a \mathrel\rho b \text{\ \ ako i samo ako $a$ i $b$ pripadaju istom klasteru.}
  $$
  Doka\v zimo da je $\rho$ relacija ekvivalencije na skupu $A$.

    (Refleksivnost)
    Neka je $a \in A$ proizvoljno. Tre\'ca osobina klasterizacije nam garantuje da se $a$ nalazi u nekom klasteru $C_i$, i zato je iskaz:
	\begin{center}
	  \enquote{$a$ i $a$ se nalaze u istom klasteru}
	\end{center}
	\v cudan, ali ta\v can. Dakle, $a \mathrel\rho a$.

    (Simetri\v cnost)
    Neka je $a \mathrel\rho b$. Tada $a$ i $b$ pripadaju istom klasteru, pa je iskaz:
	\begin{center}
	  \enquote{$b$ i $a$ nalaze u istom klasteru}
	\end{center}
	o\v cigledno ta\v can. Zato je $b \mathrel\rho a$.

    (Tranzitivnost)
    Neka je $a \mathrel\rho b$ i $b \mathrel\rho c$. Tada se $a$ i $b$ nalaze u istom klasteru, i $b$ i $c$ se nalaze u istom klasteru.
    Recimo da je $a, b \in C_i$ i $b, c \in C_j$. Kako razli\v citi klasteri nemaju zajedni\v ckih elemenata, zaklju\v cujemo
    da je $C_i = C_j$ jer $b \in C_i \cap C_j$. Dakle, $a$ i $c$ se tako\dj e nalaze u istom klasteru, pa je zato
    $a \mathrel\rho c$.

    POtrebno je jo\v s pokazati da je $A / \Boxed\rho = \{C_1, C_2, \ldots, C_k\}$, ali to sledi direktno iz konstrukcije relacije~$\rho$
	i prve osobine klasterizacije.
\end{proof}

\section{Pore\dj enje}

Sortiranje podataka po nekom kriterijumu predstavlja jedan od najstarijih i najva\v znijih procesa u obradi podataka.
Algoritam za sortiranje poredi parove podataka po nekom kriterijumu i onda odlu\v cuje koji podatak bi trebalo da se pojavi
ranije u kona\v cnom poretku.
Recimo, podatke o osobama mo\v zemo da uredimo po prezimenu, ili po starosti.
U ovom odeljku \'cemo odgovoriti na pitanje kada neka apstraktna binarna relacija $\rho \subseteq A \times A$ predstavlja
\enquote{razuman} kriterijum za pore\dj enje.

Pogledajmo jedan jednostavan primer u programskom jeziku Java. Neka nam je data klasa \texttt{Student} koja opisuje podatke o studentima, recimo ovako:
{\small
\begin{verbatim}
  class Student {
    public String ime, prezime;
    public int dan, mes, god; // datum rodjenja
    public int idx; // broj indeksa

    public Student(String i, String p, int d, int m, int g, int b) {
      ime = i; prezime = p; idx = b;
      dan = d; mes = m; god = g;
    }

    public void print() {
      System.out.println(idx + " " + ime + " " + prezime + ", "
                         + dan + "." + mes + "." + god);
    }
  }
\end{verbatim}
}
Programski jezik Java ima klasu \verb`Arrays` u paketu \texttt{java.util}
koja implementira metod \verb`sort`, ali programer mora da objasni sistemu
po kom kriterijumu ni niz trebalo da bude sortiran. Drugim re\v cima, programer mora da opi\v se relaciju poretka
koju \v zeli da primeni kako bi metod \verb`Arrays.sort` mogao da sortira elemente niza.

Apstrakcija relacije poretka u Javi je realizovana kao parametrizovani interfejs \verb`Comparator`, gde
parametar interfejsa predstavlja tip objekata koje \'cemo porediti. Ovaj interfejs sadr\v zi metod
\texttt{compare}
{\small
\begin{verbatim}
    int compare(T obj1, T obj2)
\end{verbatim}
}
\noindent
koji vra\'ca $-1$ ukoliko je \verb`obj1` strogo manje od \verb`obj2`,
vra\'ca $0$ ukoliko su \verb`obj1` i \verb`obj2` jednaki u odnosu na posmatrano svojstvo,
i vra\'ca $1$ ukoliko je \verb`obj1` strogo ve\'ci od \verb`obj2`.

Pretpostavimo da smo u glavnom programu napravili niz \texttt{S} i u njega smestili podatke
o nekoliko studenata:
{\small
\begin{verbatim}
  import java.util.*;

  class sortiranje {
    public static void main(String[] args) {
      Student[] S = {
        new Student("Petar",  "Petrovic",    1,  1, 2001, 301456),
        new Student("Jovana", "Jovanovic",   3,  2, 2002, 456701),
        new Student("Milan",  "Milankovic",  5,  7, 2002, 203467),
        new Student("Marija", "Marjanovic", 11, 12, 2001, 523801)
      };
\end{verbatim}
}
Ako \v zelimo da sortiramo studente po broju indeksa, prvo \'cemo kriterijum za sortiranje
implementirati kao instancu klase \texttt{Comparator}, a onda \'cemo iz klase \texttt{Arrays}
pozvati metod \texttt{sort} koji \'ce niz \texttt{S} sortirati po tom kriterijumu:
{\small
\begin{verbatim}
      Comparator poIdx = new Comparator<Student>() {
        @Override public int compare(Student a, Student b) {
          if(a.idx < b.idx) return -1;
          if(a.idx > b.idx) return  1;
          return 0;
        }
      };

      Arrays.sort(S, poIdx);
      for(Student a : S) a.print();
\end{verbatim}
}
Isti niz sada mo\v zemo presortirati po prezimenu, u rastu\'cem i u opadaju\'cem poretku:
{\small
\begin{verbatim}
      Comparator poPrezimenu = new Comparator<Student>() {
        @Override public int compare(Student a, Student b) {
            return a.prezime.compareTo(b.prezime);
        }
      };

      Arrays.sort(S, poPrezimenu);
      for(Student a : S) a.print();
    
      Arrays.sort(S, poPrezimenu.reversed());
      for(Student a : S) a.print();
\end{verbatim}
}
Primetimo da svaka instanca klase \texttt{Comparator} ima metod \texttt{reversed()}
koji omogu\'cuje sortiranje u obrnutom poretku od onog koji je definisan komparatorom.
Kona\v cno, ako \v zelimo studente da sortiramo po datumu ro\dj enja, definisa\'cemo odgovaraju\'ci
komparator i presortirati niz:
{\small
\begin{verbatim}
      Comparator poDatRodj = new Comparator<Student>() {
        @Override public int compare(Student a, Student b) {
            if(a.god < b.god) return -1;
            if(a.god > b.god) return 1;
            if(a.mes < b.mes) return -1;
            if(a.mes > b.mes) return 1;
            if(a.dan < b.dan) return -1;
            if(a.dan > b.dan) return 1;
            return 0;
        }
      };
    
      Arrays.sort(S, poDatRodj);
      for(Student a : S) a.print();
\end{verbatim}
}
Iako programski jezik omogu\'cuje da metod \texttt{compare} klase \texttt{Comparator} bude proizvoljan
Java metod, kriterijum za pore\dj enje mora da ispunjava neka o\v cekivanja da bismo dobili smislen rezultat.
Do apstraktnog pojma relacije poretka (tj.\ \enquote{razumnog} kriterijuma za pore\dj enje)
dolazimo analizom fundamentalnih svojstava dva prirodna poretka:
poretka na brojevima ($\le$) i poretka na skupovima ($\subseteq$).

Poredak $\le$ na brojevima ima slede\'ce osobine, gde su $x$, $y$, $z$ proizvoljni brojevi:
\begin{itemize}\itemsep -1.5pt
\item $x \le x$,
\item $x \le y \land y \le x \Rightarrow x = y$, i
\item $x \le y \land y \le z \Rightarrow x \le z$.
\end{itemize}
Interesantno je da i relacija $\subseteq$ na skupovima ima iste ove tri osobine, gde su $X$, $Y$, $Z$ sada neki skupovi:
\begin{itemize}\itemsep -1.5pt
\item $X \subseteq X$,
\item $X \subseteq Y \land Y \subseteq X \Rightarrow X = Y$, i
\item $X \subseteq Y \land Y \subseteq Z \Rightarrow X \subseteq Z$.
\end{itemize}
Ove tri osobine karakteri\v su poredak objekata, a njihovom
apstrakcijom dobijamo pojam apstraktne relacije poretka.

\begin{DEF}
  \emph{Relacija poretka na skupu $A$} je svaka binarna relacija $\Boxed\rho \subseteq A \times A$
  koja ima slede\'ce tri osobine:
  \begin{itemize}\itemsep -1pt
  \item za sve $a \in A$: $a \mathrel{\rho} a$, \hfill (\emph{refleksivnost})
  \item za sve $a, b \in A$: $a \mathrel{\rho} b \land b \mathrel{\rho} a \Rightarrow a = b$, \hfill (\emph{antisimetri\v cnost})
  \item za sve $a, b, c \in A$: $a \mathrel{\rho} b \land b \mathrel{\rho} c \Rightarrow a \mathrel{\rho} c$. \hfill (\emph{tranzitivnost})
  \end{itemize}
\end{DEF}

Neka je $\le$ uobi\v cajeni poredak celih brojeva. Ova relacija poretka osim tri navedene ima jo\v s jednu
interesantnu osobinu: \emph{svaka dva cela broja su uporediva}.
Naime, za svaka dva broja $a, b \in \ZZ$ je $a \le b$ ili je $b \le a$.
S druge strane, poredak $\subseteq$ na skupovima nema tu osobinu: postoje skupovi koji su \emph{neuporedivi}.
Na primer $\{1, 2\} \not\subseteq \{2, 3\}$ i $\{2, 3\} \not\subseteq \{1, 2\}$.

\begin{DEF}
Relacija poretka na skupu $A$ je \emph{linearni (ili totalni) poredak} ako va\v zi slede\'ce
\begin{itemize}
  \item za sve $a, b \in A$: $a \mathrel\rho b \lor b \mathrel\rho a.$\hfill (linearnost)
\end{itemize}
\end{DEF}

\begin{EX}\label{ft07.ex.deljivost}
  Neka je $\mid$ relacija deljivosti na skupu $\NN$. Na primer, $3 \mid 12$ i $5 \nmid 12$. Lako se dokazuje da je
  $\mid$ relacija poretka na $\NN$. Refleksivnost i tranzitivnost su o\v cigledni. Doka\v zimo antisimetri\v cnost.
  Neka $a \mid b$ i $b \mid a$ za neke $a, b \in \NN$. Iz $a \mid b$ sledi $a \le b$. Isto tako, iz $b \mid a$ sledi
  $b \le a$. Dakle, $a \le b$ i $b \le a$, pa je $a = b$. Ovo nije linearni poredak jer postoje neuporedivi elementi,
  recimo, $2 \nmid 3$ i $3 \nmid 2$.
\end{EX}

% \begin{EX}
%   Neka je na skupu $\NN \times \NN$ definisana relacija $\sqsubseteq$ ovako:
%   $$
%     (a, b) \sqsubseteq (c, d) \text{\ \ ako i samo ako\ \ } a \le c \land b \le d.
%   $$
%   Tada se lako vidi da je $\sqsubseteq$ relacija poretka. Ovo nije linearni poredak jer postoje neuporedivi elementi,
%   recimo, $(1, 2) \not\sqsubseteq (2, 1)$ i $(2, 1) \not\sqsubseteq (1, 2)$.
% \end{EX}

% Prethodni primer je specijalni slu\v caj slede\'ceg rezultata koji pokazuje kako mo\v ze da se konstrui\v se
% relacija poretka na proizvodu.

% \begin{THM}
%   Neka je $\sqsubseteq_1$ relacija poretka na skupu $A_1$, a $\sqsubseteq_2$ relacija poretka na skupu $A_2$.
%   Na skupu $A_1 \times A_2$ defini\v semo binarnu relaciju $\preccurlyeq$ ovako:
%   $$
%     (a_1, a_2) \preccurlyeq (b_1, b_2) \text{\ \ ako i samo ako je\ \ } a_1 \mathrel{\sqsubseteq_1} b_1 \land a_2 \mathrel{\sqsubseteq_2} b_2.
%   $$
%   Tada je $\preccurlyeq$ relacija poretka na $A_1 \times A_2$.
% \end{THM}
% \begin{proof}[Dokaz]
%   Doka\v zimo da je relacija $\preccurlyeq$ refleksivna, antisimetri\v cna i tranzitivna. Primetimo, prvo, da je
%   $$
%     (a_1, a_2) \preccurlyeq (a_1, a_2) \Leftrightarrow a_1 \mathrel{\sqsubseteq_1} a_1 \land a_2 \mathrel{\sqsubseteq_2} a_2.
%   $$
%   Kako je iskaz na desnoj strani ta\v can (zato \v sto su relacije $\sqsubseteq_1$ i $\sqsubseteq_2$ refleksivne)
%   pokazali smo da je relacija $\preccurlyeq$ refleksivna. Da bismo pokazali da je relacija $\preccurlyeq$ antisimetri\v cna
%   pretpostavimo da je $(a_1, a_2) \preccurlyeq (b_1, b_2)$ i $(b_1, b_2) \preccurlyeq (a_1, a_2)$. Prema definiciji je:
%   \begin{align*}
%     (a_1, a_2) &\preccurlyeq (b_1, b_2) \land (b_1, b_2) \preccurlyeq (a_1, a_2)\\
%                &\Leftrightarrow a_1 \mathrel{\sqsubseteq_1} b_1 \land a_2 \mathrel{\sqsubseteq_2} b_2 \land b_1 \mathrel{\sqsubseteq_1} a_1 \land b_2 \mathrel{\sqsubseteq_2} a_2\\
%                &\Leftrightarrow (a_1 \mathrel{\sqsubseteq_1} b_1 \land b_1 \mathrel{\sqsubseteq_1} a_1) \land (a_2 \mathrel{\sqsubseteq_2} b_2 \land b_2 \mathrel{\sqsubseteq_2} a_2)\\
%                &\Leftrightarrow a_1 = b_1 \land a_2 = b_2\\
%                &\Leftrightarrow (a_1, a_2) = (b_1, b_2),
%   \end{align*}
%   zato \v sto su relacije $\sqsubseteq_1$ i $\sqsubseteq_2$ antisimetri\v cne. Kona\v cno da bismo pokazali
%   da je relacija $\preccurlyeq$ tranzitivna pretpostavimo da je $(a_1, a_2) \preccurlyeq (b_1, b_2)$ i $(b_1, b_2) \preccurlyeq (c_1, c_2)$.
%   Tada je:
%   \begin{align*}
%     (a_1, a_2) &\preccurlyeq (b_1, b_2) \land (b_1, b_2) \preccurlyeq (c_1, c_2)\\
%                &\Leftrightarrow a_1 \mathrel{\sqsubseteq_1} b_1 \land a_2 \mathrel{\sqsubseteq_2} b_2 \land b_1 \mathrel{\sqsubseteq_1} c_1 \land b_2 \mathrel{\sqsubseteq_2} c_2\\
%                &\Leftrightarrow (a_1 \mathrel{\sqsubseteq_1} b_1 \land b_1 \mathrel{\sqsubseteq_1} c_1) \land (a_2 \mathrel{\sqsubseteq_2} b_2 \land b_2 \mathrel{\sqsubseteq_2} c_2)\\
%                &\Leftrightarrow a_1 \mathrel{\sqsubseteq_1} c_1 \land a_2 \mathrel{\sqsubseteq_2} c_2\\
%                &\Leftrightarrow (a_1, a_2) \preccurlyeq (c_1, c_2),
%   \end{align*}
%   zato \v sto su relacije $\sqsubseteq_1$ i $\sqsubseteq_2$ tranzitivne.
% \end{proof}


\emph{Dvokriterijumsko sortiranje} predstavlja sortiranje po dva kriterijuma. Na primer,
spisak studenata nakon prijemnog ispita se \v cesto sortira
prvo po broju osvojenih poena, dok se studenti sa istim brojem osvojenih poena unutar te grupe
sortiraju po prezimenu. Slede\'ca teorema pokazuje da dvokriterijumsko sortiranje uvek daje linearni poredak.

\begin{THM}\label{ft07.thm.lex-2}
  Neka je $\sqsubseteq_1$ linearni poredak na skupu $A_1$, a $\sqsubseteq_2$ linearni poredak na skupu $A_2$.
  Na skupu $A_1 \times A_2$ defini\v semo binarnu relaciju $\preccurlyeq$ ovako:
  $$
    (a_1, a_2) \preccurlyeq (b_1, b_2) \text{\ \ ako i samo ako je\ \ }
    (a_1 \ne b_1 \land a_1 \mathrel{\sqsubseteq_1} b_1) \lor (a_1 = b_1 \land a_2 \mathrel{\sqsubseteq_2} b_2).
  $$
  Tada je $\preccurlyeq$ linearni poredak na $A_1 \times A_2$.
\end{THM}
\begin{proof}[Dokaz]
  Doka\v zimo, prvo, da je $\preccurlyeq$ relacija poretka na $A_1 \times A_2$. Refleksivnost se pokazuje lako:
  $(a_1, a_2) \preccurlyeq (a_1, a_2)$ je ta\v cno zato \v sto je drugi disjunkt u iskazu sa desne strane ta\v can
  ($a_1 = a_1 \land a_2 \mathrel{\sqsubseteq_2} a_2$).
  
  \medskip
  
  Doka\v zimo sada antisimetri\v cnost. Neka je $(a_1, a_2) \preccurlyeq (b_1, b_2)$ i $(b_1, b_2) \preccurlyeq (a_1, a_2)$.
  Kao prvi korak dokaza\'cemo da je $a_1 = b_1$. Pretpostavimo da to nije ta\v cno. Tada je $a_1 \ne b_1$, pa na osnovu
  definicije dobijamo
  \begin{align*}
    (a_1, a_2) \preccurlyeq (b_1, b_2) &\Rightarrow a_1 \ne b_1 \land a_1 \mathrel{\sqsubseteq_1} b_1 \Rightarrow a_1 \mathrel{\sqsubseteq_1} b_1\\
    (b_1, b_2) \preccurlyeq (a_1, a_2) &\Rightarrow b_1 \ne a_1 \land b_1 \mathrel{\sqsubseteq_1} a_1 \Rightarrow b_1 \mathrel{\sqsubseteq_1} a_1,
  \end{align*}
  (zato \v sto je drugi disjunkt u iskazu sa desne strane neta\v can u oba slu\v caja). Zbog antisimetri\v cnosti relacije $\sqsubseteq_1$
  zaklju\v cujemo da je $a_1 = b_1$, \v sto je u suprotnosti sa pretpostavkom $a_1 \ne b_1$.
  
  Dakle, mora biti $a_1 = b_1$. Sada se lako pokazuje da je $a_2 = b_2$:
  \begin{align*}
    (a_1, a_2) \preccurlyeq (b_1, b_2) &\Rightarrow a_1 = b_1 \land a_2 \mathrel{\sqsubseteq_2} b_2 \Rightarrow a_2 \mathrel{\sqsubseteq_2} b_2\\
    (b_1, b_2) \preccurlyeq (a_1, a_2) &\Rightarrow b_1 = a_1 \land b_2 \mathrel{\sqsubseteq_2} a_2 \Rightarrow b_2 \mathrel{\sqsubseteq_2} a_2,
  \end{align*}
  pa je $a_2 = b_2$. Dakle, $(a_1, a_2) = (b_1, b_2)$.
  
  \medskip
  
  Da bismo pokazali tranzitivnost pretpostavimo da je $(a_1, a_2) \preccurlyeq (b_1, b_2)$ i $(b_1, b_2) \preccurlyeq (c_1, c_2)$.
  Razmotri\'cemo \v cetiri slu\v caja.
  
  \medskip
  
  Slu\v caj 1: $a_1 = b_1$ i $b_1 = c_1$. Tada je $a_2 \mathrel{\sqsubseteq_2} b_2$ i $b_2 \mathrel{\sqsubseteq_2} c_2$.
  Dakle, $a_1 = c_1$ i $a_2 \mathrel{\sqsubseteq_2} c_2$, pa je $(a_1, a_2) \preccurlyeq (c_1, c_2)$.
  
  \medskip

  Slu\v caj 2: $a_1 = b_1$ i $b_1 \ne c_1$. Tada je $b_1 \mathrel{\sqsubseteq_1} c_1$.
  Dakle, $a_1 \ne c_1$ i $a_1 \mathrel{\sqsubseteq_1} c_1$ (jer je $a_1 = b_1$), pa je $(a_1, a_2) \preccurlyeq (c_1, c_2)$.

  \medskip

  Slu\v caj 3: $a_1 \ne b_1$ i $b_1 = c_1$. Tada je $a_1 \mathrel{\sqsubseteq_1} b_1$.
  Dakle, $a_1 \ne c_1$ i $a_1 \mathrel{\sqsubseteq_1} c_1$ (jer je $b_1 = c_1$), pa je $(a_1, a_2) \preccurlyeq (c_1, c_2)$.

  \medskip

  Slu\v caj 4: $a_1 \ne b_1$ i $b_1 \ne c_1$. Tada je $a_1 \mathrel{\sqsubseteq_1} b_1$ i $b_1 \mathrel{\sqsubseteq_1} c_1$.
  Tada je $a_1 \mathrel{\sqsubseteq_1} c_1$. Lako se pokazuje da je $a_1 \ne c_1$
  (jer u suprotnom iz $a_1 = c_1$ sledi $c_1 \mathrel{\sqsubseteq_1} a_1$, pa $a_1 \mathrel{\sqsubseteq_1} b_1$ dalje implicira $c_1 \mathrel{\sqsubseteq_1} b_1$;
  to sa zaklju\v ckom $b_1 \mathrel{\sqsubseteq_1} c_1$ daje $b_1 = c_1$, \v sto je suprotno pretpostavci $b_1 \ne c_1$).
  Dakle, $a_1 \ne c_1$ i $a_1 \mathrel{\sqsubseteq_1} c_1$, pa je $(a_1, a_2) \preccurlyeq (c_1, c_2)$.
  
  \medskip
  
  Kona\v cno, doka\v zimo da je $\preccurlyeq$ linearni poredak, odnosno, da su svaka dva elementa skupa $A_1 \times A_2$
  uporediva. Neka su $(a_1, a_2), (b_1, b_2) \in A_1 \times A_2$ proizvoljni. Razmotri\'cemo dva slu\v caja.
  
  \medskip
  
  Slu\v caj 1: $a_1 = b_1$. Kako je $\sqsubseteq_2$ linearni poredak, mora bili $a_2 \mathrel{\sqsubseteq_2} b_2$
  ili $b_2 \mathrel{\sqsubseteq_2} a_2$, pa je, u zavisnosti od toga koja od ove dve relacije je ta\v cna,
  $(a_1, a_2) \preccurlyeq (b_1, b_2)$ ili $(b_1, b_2) \preccurlyeq (a_1, a_2)$.
  
  \medskip
  
  Slu\v caj 2: $a_1 \ne b_1$. Kako je $\sqsubseteq_1$ linearni poredak, mora bili $a_1 \mathrel{\sqsubseteq_1} b_1$
  ili $b_1 \mathrel{\sqsubseteq_1} a_1$, pa je, u zavisnosti od toga koja od ove dve relacije je ta\v cna,
  $(a_1, a_2) \preccurlyeq (b_1, b_2)$ ili $(b_1, b_2) \preccurlyeq (a_1, a_2)$.
  
  \medskip
  
  Ovim je tvr\dj enje pokazano.
\end{proof}


\section{Pore\dj enje re\v ci}

Da se podsetimo, za kona\v can skup $A$ jo\v s ka\v zemo i da je \emph{alfabet}.
Za njegove elemente tada ka\v zemo da su \emph{slova}.
Ranije smo videli da je re\v c nad alfabetom $A$ proizvoljan kona\v can niz slova tog alfabeta,
uklju\v cuju\'ci i praznu re\v c koju smo ozna\v cili sa~$\epsilon$.
Skup svih re\v ci nad alfabetom $A$ je, dakle, skup svih kona\v cnih nizova elemenata skupa~$A$:
$$
  A^* = \{\epsilon\} \cup A^1 \cup A^2 \cup \ldots \cup A^n \cup \ldots.
$$
Re\v c $u$ je \emph{prefiks} re\v ci $v$ ako se re\v c $u$ nalazi na samom po\v cetku re\v ci $v$, iza \v cega u re\v ci
$v$ mo\v ze, a ne mora da se nalazi jo\v s slova. Relacija \enquote{\ldots\ je prefiks od \ldots} je relacija poretka
na $A^*$ koja ne mora biti linearni poredak. Na primer, za $A = \{a, b, c\}$ imamo slede\'ce:
\begin{itemize}\itemsep -1.5pt
\item
  $ab$ je prefiks od $abc$, a nije prefiks od $aab$;
\item
  prazna re\v c je prefiks svake re\v ci;
\item
  za svaku re\v c $u$ je $u$ prefiks od $u$ (refleksivnost);
\item
  ako je $u$ prefiks od $v$ i ako je $v$ prefiks od $u$, onda je $u = v$ (antisimetri\v cnost);
\item
  ako je $u$ prefiks od $v$ i ako je $v$ prefiks od $w$, onda je $u$ prefiks od $w$ (tranzitivnost);
\item
  niti je $ab$ prefiks od $ba$, niti je $ba$ prefiks od $ab$ (zato ovo nije linearni poredak).
\end{itemize}

\emph{Leksikografski poredak} predstavlja na\v cin da se re\v ci nekog jezika urede \emph{linearno.}
Zove se \emph{leksikografski} jer se isti koristi za ure\dj ivanje re\v ci u leksikonu.
Recimo:
\begin{quote}
  Avala,
  Besna kobila,
  Cer,
  Dukat,
  Fruška gora,
  Golija,
  Homoljske planine,
  Jelova gora,
  Kopaonik,
  Lepa gora,
  Maljen,
  Ozren,
  Povlen,
  Rtanj,
  Stara planina,
  Tara,
  Vršačka brda,
  Zlatibor
\end{quote}
Va\v zno je uo\v citi da leksikografski poredak re\v ci zavisi od poretka slova u alfabetu.

Neka je $A$ alfabet i neka je $\sqsubseteq$ proizvoljan linearni poredak na $A$. Na skupu $A^*$
defini\v semo linearni poredak $\preccurlyeq$ ovako. Neka su $u = a_1 a_2 \ldots a_m$ i $v = b_1 b_2 \ldots b_n$ proizvoljne
re\v ci iz~$A^*$.
\begin{itemize}
\item
    Ako je $u$ prefiks od $v$ onda je $u \preccurlyeq v$; ako je $v$ prefiks od $u$ onda je $v \preccurlyeq u$.
\item
    Ako nijedna od ove dve re\v ci nije prefiks one druge uo\v cimo najmanje $i$ takvo da je $a_i \ne b_i$,
    (to zna\v ci da je $a_1 = b_1$, $a_2 = b_2$, \ldots, $a_{i-1} = b_{i-1}$), pa za $a_i \sqsubseteq b_i$ stavimo
    $u \preccurlyeq v$, dok za $b_i \sqsubseteq a_i$ stavimo $v \preccurlyeq u$.
\end{itemize}
Tada je $\preccurlyeq$ linearni poredak na~$A^*$.
Na primer, za $A = \{a, b, c, d\}$ imamo slede\'ce: $aaaaa \preccurlyeq b \preccurlyeq bcac \preccurlyeq bcd$.



\section{Grafi\v cki prikaz relacija poretka}


Relacije poretka kao skupovi ure\dj enih parova obi\v cno imaju mnogo elemenata, pa grafi\v cko predstavljanje koje smo
koristili kod relacija ekvivalencije nije pogodno. Istini za volju, nije bilo pogodno ni za relacije ekvivalencije --
ve\'c mali primeri su bili predstavljeni \v sumama linija, pa smo smislili efikasniji na\v cin (koli\v cni\v cki skup)
da predstavimo relaciju ekvivalencije.

Relacije poretka se grafi\v cki predstavljaju koriste\'ci zgodnu dosetku koja se zove Haseov dijagram.
\emph{Haseov dijagram} prikazuje relaciju poretka tako \v sto na dijagram unosi minimalnu
koli\v cinu informacija, uz slede\'ce dogovore:
\begin{itemize}
  \item ako je element $a$ manji od elementa $b$, nacrta\'cemo $a$ negde ispod $b$ i spoji\'cemo ih linijom;
        linija nema strelicu jer se podrazumeva da je orijentisana od elementa koji je ni\v ze na crte\v zu
        ka elementu koji je nacrtan iznad njega;
  \item ne crtamo petlje (znamo da je relacija refleksivna);
  \item ne crtamo linije koje mo\v zemo da dobijemo kao posledicu nacrtanih linija i \v cinjenice da je relacija tranzitivna.
\end{itemize}
Ne\'cemo ovde ulaziti u formalnu raspravu o tome kako se ta\v cno defini\v se Haseov dijagram,
ve\'c \'cemo na nekoliko primera pokazati kako se on konstrui\v se za datu relaciju poretka.

\begin{EX}\label{ft07.ex.deljivost-1-10}
  Relacija deljivosti $\mid$ je relacija poretka na skupu $\{1, 2, 3, 4, 5, 6, 7, 8, 9, 10\}$
  (Primer~\ref{ft07.ex.deljivost}). Pokaza\'cemo sada kako se crta Haseov dijagram koji joj odgovara.
\end{EX}

\bigskip

\noindent
\begin{minipage}{3in}
  Crtanje kre\'cemo od nekog zgodnog elementa, \v sto je u ovom slu\v caju~1.
\end{minipage}
\hfill
\begin{minipage}{2.5in}
  \begin{center}
    \input ft07-r6a.pgf
  \end{center}
\end{minipage}

\bigskip

\noindent
\begin{minipage}{3in}
  Slede\'ci na redu je element~2. Kako $1 \mid 2$, nacrta\'cemo~2 iznad~1 i spojiti ih linijom.
\end{minipage}
\hfill
\begin{minipage}{2.5in}
  \begin{center}
    \input ft07-r6b.pgf
  \end{center}
\end{minipage}

\bigskip

\noindent
\begin{minipage}{3in}
  Broj 3 je deljiv sa 1, ali nije deljiv sa 2. Zato \'cemo broj 3 nacrtati iznad 1 i spojiti ga linijom sa 1, a ne\'cemo
  ga spojiti linijom sa 2.
\end{minipage}
\hfill
\begin{minipage}{2.5in}
  \begin{center}
    \input ft07-r6c.pgf
  \end{center}
\end{minipage}

\bigskip

\noindent
\begin{minipage}{3in}
  Broj 4 je deljiv brojem 2, ali nije deljiv brojem 3. Zato \'cemo broj 4 nacrtati iznad broja 2 i spojiti ga linijom sa 2, a ne\'cemo
  ga spojiti linijom sa 3.
\end{minipage}
\hfill
\begin{minipage}{2.5in}
  \begin{center}
    \input ft07-r6d.pgf
  \end{center}
\end{minipage}

\medskip

Primetimo da je broj 4 deljiv i brojem 1. Me\dj utim, ne\'cemo broj 4 spojiti linijom sa brojem 1 da bismo smanjili broj linija na
dijagramu. Zato se kod Haseovih dijagrama uvodi slede\'ci dogovor:
\begin{quote}
    element $x$ je manji od elementa $y$ ako je $x$ ispod $y$ i postoji
    put koji spaja $x$ sa $y$; ako takav put ne postoji, onda su $x$ i $y$ neuporedivi ili je $y$ manji of $x$.
\end{quote}
To je slu\v caj sa brojevima 3 i 4 na prethodnom dijagramu: kako su 3 i 4 neuporedivi,
ne postoji put od 3 do 4 koji ide samo navi\v se, niti put od 4 do 3 koji ide samo navi\v se.

\bigskip

\noindent
\begin{minipage}{3in}
  Od svih navedenih brojeva broj 5 je deljiv samo brojem 1. Zato \'cemo broj 5 nacrtati iznad broja 1 i spojiti ga linijom sa 1, a ne\'cemo
  ga spojiti linijom ni sa jednim drugim elementom na slici.
\end{minipage}
\hfill
\begin{minipage}{2.5in}
  \begin{center}
    \input ft07-r6e.pgf
  \end{center}
\end{minipage}

\bigskip

\noindent
\begin{minipage}{3in}
  Broj 6 je deljiv brojevima 2 i 3. Zato \'cemo broj 6 nacrtati iznad brojeva 2 i 3 i spojiti ga linijom i sa 2 i sa 3.
  Broj 6 je deljiv i brojem 1, ali smo se dogovorili da ne crtamo linije ukoliko postoji put koji spaja uporedive elemente i kre\'ce
  se samo navi\v se.
\end{minipage}
\hfill
\begin{minipage}{2.5in}
  \begin{center}
    \input ft07-r6f.pgf
  \end{center}
\end{minipage}

\bigskip

\noindent
\begin{minipage}{3in}
  Od svih navedenih brojeva broj 7 je deljiv samo brojem 1. Zato \'cemo broj 7 nacrtati iznad broja 1 i spojiti ga linijom sa 1, a ne\'cemo
  ga spojiti linijom ni sa jednim drugim elementom na slici.
\end{minipage}
\hfill
\begin{minipage}{2.5in}
  \begin{center}
    \input ft07-r6g.pgf
  \end{center}
\end{minipage}

\bigskip

\noindent
\begin{minipage}{3in}
  Od svih navedenih brojeva broj 8 je deljiv samo brojevima 1, 2 i 4. Zato \'cemo broj 8 nacrtati iznad broja 4 i spojiti ga linijom
  samo sa brojem 4. To je, prema dogovoru, dovoljno da se nazna\v ci da je 8 deljiv i sa 4, i sa 2, i sa 1.
\end{minipage}
\hfill
\begin{minipage}{2.5in}
  \begin{center}
    \input ft07-r6h.pgf
  \end{center}
\end{minipage}

\bigskip

\noindent
\begin{minipage}{3in}
  Sli\v cno tome, broj 9 je deljiv brojevima 1 i 3, pa \'cemo broj 9 nacrtati iznad broja 3 i spojiti ga linijom
  samo sa brojem 3.
\end{minipage}
\hfill
\begin{minipage}{2.5in}
  \begin{center}
    \input ft07-r6i.pgf
  \end{center}
\end{minipage}

\bigskip

\noindent
\begin{minipage}{3in}
  Kona\v cno, broj 10 je deljiv brojevima 2 i 5, pa \'cemo ga nacrtati iznad brojeva 2 i 5, i spojiti ga linijom
  samo sa oba. Time je Haseov dijagram u ovom primeru kompletiran.
\end{minipage}
\hfill
\begin{minipage}{2.5in}
  \begin{center}
    \input ft07-r6z.pgf
  \end{center}
\end{minipage}

\bigskip

\noindent
\begin{minipage}{3in}
  \begin{EX}\label{ft07.ex.P123}
    Pored je dat Haseov dijagram relacije poretka $\subseteq$ na skupu
    $\calP(\{a, b, c\})$.
  \end{EX}
\end{minipage}
\hfill
\begin{minipage}{2.5in}
  \begin{center}
    \input ft07-r7.pgf
  \end{center}
\end{minipage}

\bigskip

\noindent
\begin{minipage}{3in}
  \begin{EX}
    Pored je dat Haseov dijagram relacije poretka $\le$ na skupu
    $\{1, 2, 3, 4, 5\}$. Svaki linearni poredak ima ovakav Haseov dijagram.
    Ime \emph{linearni poredak} poti\v ce od \v cinjenice da je Haseov dijagram
    takvog ure\dj enja \enquote{jedna linija}.
  \end{EX}
\end{minipage}
\hfill
\begin{minipage}{2.5in}
  \begin{center}
    \input ft07-r8.pgf
  \end{center}
\end{minipage}

\bigskip

\emph{Najmanji element} relacije poretka na skupu $A$ je element skupa $A$ koji je uporediv sa svim elementima skupa $A$
i nalazi se ispod svih njih. Sli\v cno, \emph{najve\'ci element} relacije poretka na skupu $A$ je element skupa $A$
koji je uporediv sa svim elementima skupa $A$ i nalazi se iznad svih njih. Preciznije, neka je $\sqsubseteq$ relacija
poretka na skupu~$A$. Element $a \in A$ je \emph{najmanji} element ako
$$
  (\forall x \in A) \; a \sqsubseteq x,
$$
a element $b \in A$ je \emph{najve\'ci} element ako
$$
  (\forall x \in A) \; x \sqsubseteq b.
$$

\begin{EX}\label{ft07.ex.delj-2-12}
  Uobi\v cajena relacija poretka $\le$ na skupu $\ZZ$ nema ni najmanji ni najve\'ci element.

  Relacija poretka koja je prikazana u Primeru~\ref{ft07.ex.deljivost-1-10} ima najmanji element, ali nema najve\'ci.

  Relacija poretka koja je prikazana u Primeru~\ref{ft07.ex.P123} ima i najmanji i najve\'ci element.
\end{EX}

\bigskip

\noindent
\begin{minipage}{3in}
  \hspace*{1.5em}Relacija poretka $|$ (deljivost) na skupu $\{2, 3, 4, 6, 12\}$ je prikazana pored. Ona ima najve\'ci element, ali
  nema najmanji element.
\end{minipage}
\hfill
\begin{minipage}{2.5in}
  \begin{center}
    \input ft07-r10.pgf
  \end{center}
\end{minipage}

\bigskip

Pogledajmo jo\v s jednom relacije poretka iz Primera~\ref{ft07.ex.deljivost-1-10} i~\ref{ft07.ex.delj-2-12}:
\begin{center}
    \input ft07-r6z.pgf
    \hspace*{2cm}
    \input ft07-r10.pgf
\end{center}
Vidimo da prva relacija ima najmanji element, ali nema najve\'ci element, dok druga relacija ima najve\'ci element, ali nema
najmanji element. Ako pogledamo element 10 u prvoj relaciji, vidimo da je on nekako \enquote{pri vrhu} i da nije najve\'ci element samo
zato \v sto nije uporediv sa elementima kao \v sto su 6, 7, 8, 9. Za takav element relacije poretka ka\v zemo da je
\emph{maksimalan element}. Sli\v cno, element 2 druge relacije poretka je nekako \enquote{pri dnu} i to nije najmanji element
samo zato \v sto nije uporediv sa~3. Za takav element relacije poretka ka\v zemo da je \emph{minimalan element}.

Preciznije, neka je $\sqsubseteq$ relacija poretka na skupu $A$. Element $a \in A$ je \emph{maksimalan element} ako
$$
  \lnot(\exists x \in A)(a \ne x \land a \sqsubseteq x),
$$
odnosno, ako iznad njega nema drugih elemenata. Sli\v cno, element $b \in A$ je \emph{minimalan element} ako
$$
  \lnot(\exists x \in A)(x \ne b \land x \sqsubseteq b),
$$
odnosno, ako ispod njega nema drugih elemenata. Relacija poretka mo\v ze, a ne mora imati maksimalne i minimalne elemente,
minimalnih i maksimalnih elemenata mo\v ze biti vi\v se, a mo\v ze se desiti da je neki element istovremeno i maksimalan
i minimalan element relacije pretka, kao \v sto pokazuju slede\'ci primeri.

\begin{EX}
  Relacija poretka $\mid$ (deljivost) na skupu $\NN$ ima najmanji element~1, a nema maksimalne elemente zato \v sto za svaki
  prirodan broj $m$ postoji prirodan broj $n > m$ takav da je $m \mid n$.
\end{EX}

\bigskip

\noindent
\begin{minipage}{3in}
  \begin{EX}
    Relacija poretka $|$ (deljivost) na skupu $\{2, 3, 4, 5, 6\}$ je prikazana pored. Minimalni elementi ove relacije
    su 2, 3 i 5, jer ispod njih nema drugih elemenata. Maksimalni elementi ove relacije su 4, 6 i 5, jer iznad njih nema
    drugih elemenata. Vidimo da je 5 istovremeno i minimalni i maksimalni element.
  \end{EX}
\end{minipage}
\hfill
\begin{minipage}{2.5in}
  \begin{center}
    \input ft07-r11.pgf
  \end{center}
\end{minipage}

\bigskip

Relacija poretka ne mora da ima ni najmanji ni najve\'ci element, ali ako takvi elementi postoje onda su jedinstveni,
kako to pokazuje slede\'ca teorema.

\begin{THM}
  Neka je $\sqsubseteq$ relacija poretka na skupu $A$.
  \begin{enumerate}
  \item
    Ako relacija poretka $\sqsubseteq$ ima najve\'ci element, onda je to ujedno i njen jedinstveni maksimalni element.
  \item
    Ako relacija poretka $\sqsubseteq$ ima najmanji element, onda je to ujedno i njen jedinstveni minimalni element.
  \end{enumerate}
\end{THM}

\begin{DEF}
  Ako je na skupu $A$ uo\v cena relacija poretka $\sqsubseteq$, onda za $A$ ka\v zemo da je
  \emph{parcijalno ure\dj eni skup} i ozna\v cavamo ga kao par $(A, \sqsubseteq)$.
  
  \emph{Topolo\v sko sortiranje} je proces kojim se elementi parcijalno ure\dj enog skupa $(A, \sqsubseteq)$
  pore\dj aju u niz:
  $$
    a_0, a_1, \ldots, a_{n-1}
  $$
  tako da va\v zi slede\'ce: ako je $x \sqsubseteq y$, onda se $x$ javlja u nizu pre $y$. 
\end{DEF}

\begin{EX}
Topolo\v skim sortiranjem parcijalno ure\dj enog skupa
\begin{center}
  \input ft07-r10.pgf
\end{center}
mo\v zemo dobiti svaki od slede\'cih nizova:
$$
  \begin{array}{l}
  2, 3, 4, 6, 12\\
  3, 2, 4, 6, 12\\
  2, 3, 6, 4, 12\\
  3, 2, 6, 4, 12\\
  2, 4, 3, 6, 12
  \end{array}
$$
\emph{jer rezultat topolo\v skog sortiranja nije jednozna\v can}!
\end{EX}

\begin{EX}
  Sortiranje formiranjem uzastopnih minimuma (tzv.\ \emph{selection sort}) uz minimalne modifikacije
  mo\v ze da re\v si problem topolo\v skog sortiranja parcijalno ure\dj enih skupova.
  
  Pretpostavi\'cemo da su elementi skupa $A$ navedeni u nizu \texttt{int[] A} u proizvoljnom redosledu i
  da nam je dat metod \texttt{static boolean uporedi(int x, int y)} koji za dva elementa
  skupa $A$ vra\'ca true ako i samo ako je $x \sqsubseteq y$.

{\small
\begin{verbatim}
  static boolean uporedi(int x, int y) { ... }
  
  static void swap(int[] A, int i, int j) {
    int tmp = A[i];
    A[i] = A[j];
    A[j] = tmp;
  }

  static int posMin(int[] A, int i) {
    int n = A.length;
    int m = i;
    for(int j = i + 1; j < n; j++) {
        if(uporedi(A[j], A[m])) { m = j; }
    }
    return m;
  }

  static void TopSort(int[] A) {
    int n = A.length;
    for(int i = 0; i < n; i++) {
      int m = posMin(A, i);
      swap(A, i, m);
    }
  }
\end{verbatim}
}
\end{EX}

\clearpage

\section{Zadaci}

\ZAD{08.zad.1}{
Dokazati slede\' ce tvrdnje:
\begin{enumerate}[(a)]
\item Refleksivnost i simetri\v cnost ne impliciraju tranzitivnost.
\item Refleksivnost i tranzitivnost ne impliciraju simetri\v cnost.
\item Simetri\v cnost i tranzitivnost ne impliciraju refleksivnost.
\item Simetri\v cnost, antisimetri\v cnost i tranzitivnost ne impliciraju refleksivnost.
\item Relacija koja je simetri\v cna i antisimetri\v cna je tranzitivna.
\end{enumerate}
}

\ZAD{08.zad.2}{
Neka je $X=\{0,1,2,3,4,5\}$, i neka je $\rho$ relacija na $X$ definisana sa:
\begin{center}
$x\mathrel\rho y$ akko $x+y=5\vee x-y=0$.
\end{center}
\begin{enumerate}[(a)]
\item Nacrtati graf relacije $\rho$,
\item Odgovoriti na pitanja da li je $\rho$ refleksivna, simetri\v cna, antisimetri\v cna, odnosno tranzitivna.
\item Da li je $\rho$ relacija ekvivalencije, i da li je $\rho$ relacija poretka?
\end{enumerate}
}

\ZAD{08.zad.3}{
Neka je $X=\{6,10,11,30,35,49\}$, i neka je $\rho$ relacija na $X$ definisana sa:
\begin{center}
$x\mathrel\rho y$ akko $\NZD(x,y)\neq 1$.
\end{center}
\begin{enumerate}[(a)]
\item Nacrtati graf relacije $\rho$,
\item Odgovoriti na pitanja da li je $\rho$ refleksivna, simetri\v cna, antisimetri\v cna, odnosno tranzitivna.
\item Da li je $\rho$ relacija ekvivalencije, i da li je $\rho$ relacija poretka?
\end{enumerate}
}

\ZAD{08.zad.4}{
Neka je $X=\{0,1,2,3,4,5\}$, i neka je $\rho$ relacija na $X$ definisana sa:
\begin{center}
$x\mathrel\rho y$ akko $x-y\equiv_6 1\ \vee\  y-x\equiv_6 2\ \vee\  y=0$.
\end{center}
\begin{enumerate}[(a)]
\item Nacrtati graf relacije $\rho$,
\item Odgovoriti na pitanja da li je $\rho$ refleksivna, simetri\v cna, antisimetri\v cna, odnosno tranzitivna.
\item Da li je $\rho$ relacija ekvivalencije, i da li je $\rho$ relacija poretka?
\end{enumerate}
}

\ZAD{08.zad.5}{
Neka je $X=\{0,1,2,3,4\}$, i neka je $\rho$ relacija na $X$ definisana sa:
\begin{center}
$x\mathrel\rho y$ akko $|x-y|\equiv_5 1\text{ ili } |x-y|\equiv_5 2$.
\end{center}
\setlength{\multicolsep}{\topsep}
\begin{enumerate}[(a)]
\item Nacrtati graf relacije $\rho$,
\item Odgovoriti na pitanja da li je $\rho$ refleksivna, simetri\v cna, antisimetri\v cna, odnosno tranzitivna.
\item Da li je $\rho$ relacija ekvivalencije, i da li je $\rho$ relacija poretka?
\end{enumerate}
}

\ZAD{08.zad.6}{
Odgovoriti na pitanja da li je relacija $\rho$ nad skupom prirodnih brojeva refleksivna, simetri\v cna, antisimetri\v cna i tranzitivna, ako je:
\begin{enumerate}[(a)]
\item $x\mathrel\rho y$ akko $3\mid x-y$;
\item $x\mathrel\rho y$ akko $x=3^ky$ za neko $k\in\mathbb N$;
\item $x\mathrel\rho y$ akko $x\neq y$;
\item $x\mathrel\rho y$ akko $y\mid x$.
\end{enumerate}
}

\ZAD{08.zad.7}{
Odgovoriti na pitanja da li je relacija $\rho$ nad skupom pravih jedne ravni refleksivna, simetri\v cna, antisimetri\v cna, odnosno tranzitivna, ako je:
\begin{enumerate}[(a)]
\item $x\mathrel\rho y$ akko $x\parallel y$;
\item $x\mathrel\rho y$ akko $x\perp y$.
\end{enumerate}
}

\ZAD{08.zad.8}{
Neka je $A$ skup svih iskaznih formula. Neka je $\rho$ relacija na $A$ definisana sa
\begin{center}
$x\mathrel \rho y$ akko je $x\Rightarrow y$ tautologija.
\end{center}
Odgovoriti na pitanja da li je $\rho$ refleksivna, simetri\v cna, antisimetri\v cna, odnosno tranzitivna.
}

\begin{NAP}\it
Binarne relacije nad skupom $A$ su podskupovi skupa $A \times A$, pa se od njih mogu konstruisati
nove relacije $\rho \cup \sigma$ i $\rho \cap \sigma$ kao unija, odnosno, presek odgovaraju\'cih skupova.
Pored toga, sa $\rho\circ\sigma$ ozna\v cavamo relaciju tako\dj e nad skupom $A$ takvu da je
$x\mathrel{\rho\circ\sigma} y$ akko postoji element $z$ takav da je $x\mathrel\rho z$ i $z\mathrel\sigma y$.
\end{NAP}

\ZAD{08.zad.10}{
Ako su $\rho$ i $\sigma$ antisimetri\v cne relacije, odgovoriti na pitanje da li je $\rho\cap\sigma$,
$\rho\cup\sigma$, odnosno $\rho\circ\sigma$ antisimetri\v cna?
Odgovoriti na analogna pitanja za refleksivnost, simetri\v cnost i tranzitivnost.
}

\begin{NAP}\it
Ako je $\rho$ binarna relacija nad skupom $A$, tada sa $\rho^{-1}$ ozna\v cavamo binarnu relaciju nad istim skupom takvu da je $x\mathrel \rho^{-1}y$ akko je $y\mathrel\rho x$.
\end{NAP}

\ZAD{08.zad.9}{
Ako je $\rho$ antisimetri\v cna, da li je i $\rho^{-1}$ antisimetri\v cna? Odgovoriti na analogna pitanja za refleksivnost, simetri\v cnost i tranzitivnost.
}

\ZAD{08.zad.11}{
Proveriti da li je $\rho$ relacija ekvivalencije na ravni $\mathbb R^2$ i opisati njene klase ekvivalencija, ako je:
\begin{enumerate}[(a)]
\item $(x_1,y_1)\mathrel \rho (x_2,y_2)$ akko $x_1=x_2$;
\item $(x_1,y_1)\mathrel \rho (x_2,y_2)$ akko $x_1+y_1=x_2+y_2$;
\item $(x_1,y_1)\mathrel \rho (x_2,y_2)$ akko $x_1^2+y_1^2=x_2^2+y_2^2$.
\end{enumerate}
}\label{prethodni-zadatak}

\ZAD{08.zad.12}{
Proveriti da li je $\rho$ relacija ekvivalencije na ravni $\mathbb R^2$, gde je:
\begin{center}
$(m,n)\mathrel\rho(p,q)$ akko $m+q=n+p$.
\end{center}
Odrediti klase ekvivalencije elemenata $(1,1)$, $(2,1)$, $(3,1)$ i $(1,3)$. Opisati sve njene klase ekvivalencije.
}

\ZAD{08.zad.13}{
Koliko ima relacija ekvivalencije nad skupom:
\begin{enumerate}[(a)]
\item $\{a,b,c\}$;
\item $\{a,b,c,d\}$.
\end{enumerate}
}

\ZAD{09.zad.1}{
Da li je relacija deljivosti parcijalno ure\dj enje na skupu prirodnih brojeva? Da li ona ima najmanji, i da li ima najve\' ci element? Odgovoriti na ista pitanja za relaciju deljivosti na slupu celih brojeva.
}

\ZAD{09.zad.2}{
 Odgovoriti za\v sto $\rho$ nije parcijalno ure\dj enje na skupu $A$, ako je:
\begin{enumerate}[(a)]
\item $A=\mathcal P(\{1,2,3\})$; $B\mathrel\rho C$ akko $B\subsetneq C$;
\item $A=\mathcal P(\{1,2,3\})$; $B\mathrel\rho C$ akko $|B|\leq |C|$;
\item $A=\mathbb Z$;  $n\mathrel\rho m$ akko $n^2\leq m^2$;
\item $A=\mathbb R\times \mathbb R$; $(x_1,y_1)\mathrel\rho(x_2,y_2)$ akko $x_1\leq x_2$;
\item $A$ je grupa ljudi u kojoj nikoje dve osobe nemaju isti broj godina, niti su iste visine. Za osobe $P$ i $Q$ iz skupa $A$ pi\v semo $P\mathrel\rho Q$ akko je
\begin{center}
$\mathrm{vis}(P)\leq\mathrm{vis}(Q)\vee\mathrm{god}(P)\leq\mathrm{god}(Q)$,
\end{center}
gde $\mathrm{vis}(X)$ i $\mathrm{god}(X)$ ozna\v cavaju visinu osobe $X$ izra\v zenu u centimetrima i broj godina osobe $X$, redom.
\item $A=\mathbb Z\times\mathbb Z$; $(n,m)\mathrel\rho(p,q)$ akko $n\leq p\vee m\leq q$.
\end{enumerate}
}

\ZAD{09.zad.3}{
Ako je $\rho$ relacija poretka na $A$, da li je $\sigma$ relacija poretka na $A\times A$, ako je $(x_1,y_1)\mathrel\sigma(x_2,y_2)$ akko $x_1\mathrel\rho x_2$ i $y_1\mathrel\rho y_2$?
}

\ZAD{09.zad.4}{
Dokazati da je $\rho$ definisana sa $x\mathrel\rho y$ akko $x=2^ky$ za neko $k\in\mathbb N_0$ relacija poretka na $\mathbb N$.
}

\ZAD{09.zad.5}{
Dokazati da je $\rho$ definisana sa $x\mathrel\rho y$ akko $x=2^ky$ za neko $k\in\mathbb Z$ relacija ekvivalencije na $\mathbb N$.
}

\begin{NAP}\it
Ka\v zemo da je binarna relacija $\rho$ \emph{irefleksivna} ako ne postoji element $x$ takav da je $x\mathrel\rho x$.
Relacija \emph{strogog poretka} je binarna relacija koja je irefleksivna, antisimetri\v cna i tranzitivna.
\end{NAP}

\ZAD{09.zad.6}{
Neka je $\rho$ relacija strogog poretka na $A$. Dokazati da je $\sigma$ relacija poretka na $A$ ako je $x\mathrel\sigma y$ akko $x\mathrel\rho y\vee x=y$.
}

%\ZAD{09.zad.7}{
%Totalno ure\dj enje $\rho$ je dobro ure\dj enje na $X$ ako svaki podskup od $X$ ima najmanji element.
%\begin{enumerate}[(a)]
%\item Dokazati da je ure\dj enje na svakom kona\v cnom skupu dobro ure\dj enje.
%\item Prona\' ci primer beskona\v cnog dobro ure\dj enog skupa.
%\item Pokazati da relacija $\leq$ nije dobro ure\dj enje na skupu $\mathbb Z$.
%\end{enumerate}\footnote{Dokaz presko\v cen.}
%}

\ZAD{09.zad.8}{
Nacrtati Hase dijagram za svaki od slede\' cih skupova i za relaciju deljivosti:
\setlength{\multicolsep}{\topsep}
\begin{multicols}{2}
\begin{enumerate}[(a)]
\item $\{1,2,3,4,6,12\}$;
\item $\{2,3,4,6,8,7\}$;
\item $\{1,2,4,5,10,20\}$;
\item $\{1,2,4,8,16,32\}$;
\item $\{1,2,3,5,6,10,15,30\}$;
\item $\{2,3,5,6,10,15,30\}$.
\end{enumerate}
\end{multicols}
Za svaki od slu\v cajeva odrediti minimalne elemente, maksimalne elemente, najmanji element, najve\' ci element i lanac najve\' ce du\v zine.
}

\ZAD{09.zad.9}{
Neka je $A=\{0,1,2\}\times \{2,5,8\}$. Proveriti da li je $\rho$ parcijalno ure\dj enje na $A$ ako je
\begin{center}
$(a,b)\mathrel\rho(c,d)$ akko $a+b\mid c+d$. 
\end{center}
\begin{enumerate}[(a)]
\item Nacrtati Hase dijagram za poset $(A,\rho)$.
\item Odrediti sve maksimalne i minimalne elemente? Da li postoje najmanji i najve\' ci element?
\end{enumerate}
}

\ZAD{09.zad.10}{
Neka je $\rho$ relacija na skupu prirodnih brojeva $\mathbb N$ definisana na slede\' ci na\v cin:
\begin{center}
$x\mathrel\rho y$ akko $x=y\vee \big(x\mid y\wedge(\exists p\in\mathcal P)(p\mid y\wedge p\mathrel{\not|}x)\big),$
\end{center}
gde je $\mathcal P$ skup prostih brojeva.
\begin{enumerate}[(a)]
\item Nacrtati graf relacije $\rho$ sa brojevima 1, 2, 3, 4, 5, 6, 8, 10 i 12.
\item Odgovoriti na pitanja da li je $\rho$ refleksivna, simetri\v cna, antisimetri\v cna i tranzitivna.
\item Proveriti da li je $\rho$ relacija ekvivalencije, i da li je relacija poretka.
\item Ako je $\rho$ relacija poretka, nacrtati njen Hase dijagram sa brojevima 1, 2, 3, 4, 5, 6, 8, 9, 10, 12, 15, 18, 30 i 60, i proveriti da li je $\rho$ relacija totalnog poretka.
\end{enumerate}
}

\ZAD{09.zad.11}{
Neka je $\rho$ relacija na skupu prirodnih brojeva $\mathbb N$ definisana na slede\' ci na\v cin:
\begin{center}
$x\mathrel\rho y$ akko $x=y\vee (\exists k\in\mathbb N)(x\mid k\wedge y=kx)$.
\end{center}
\begin{enumerate}[(a)]
\item Nacrtati graf relacije $\rho$ sa brojevima 1, 2, 3, 4, 6, 8, 9 i 12.
\item Odgovoriti na pitanja da li je $\rho$ refleksivna, simetri\v cna, antisimetri\v cna i tranzitivna.
\item Proveriti da li je $\rho$ relacija ekvivalencije, i da li je relacija poretka.
\item Ako je $\rho$ relacija poretka, nacrtati njen Hase dijagram sa brojevima 2, 3, 5, 6, 8, 9, 10, 12, 16, 18, 20, 25, 30, 36 i 72, i proveriti da li je $\rho$ relacija totalnog poretka.
\end{enumerate}
}

\ZAD{09.zad.12}{
Neka je $\rho$ relacija na skupu prirodnih brojeva $\mathbb N$ definisana na slede\' ci na\v cin:
\begin{center}
$x\mathrel\rho y$ akko $(\exists n\in\mathbb N)(y=n^2x)$.
\end{center}
\begin{enumerate}[(a)]
\item Nacrtati graf relacije $\rho$ sa prvih deset prirodnih brojeva.
\item Odgovoriti na pitanja da li je $\rho$ refleksivna, simetri\v cna, antisimetri\v cna i tranzitivna.
\item Proveriti da li je $\rho$ relacija ekvivalencije, i da li je relacija poretka.
\item Ako je $\rho$ relacija poretka, nacrtati njen Hase dijagram (sa prvih 20 prirodnih brojeva), i proveriti da li je $\rho$ relacija totalnog poretka.
\end{enumerate}
}

